\begin{minipage}{0.64\linewidth}
    \begin{donnees}
        \begin{itemize}
          \item Masse de la partie immergée de l'électrode de Zinc~:
                $m=\qty{6,54}{\gram}$;
          \item Volume de chaque solution~: $V=\qty{50}{\mL}$;
          \item Concentration de chaque solution~: $C=\qty{1,0}{\M}$;
          \item $1~\mathcal{F}=\qty{9,65e4}{\coulomb\per\mole}$;
          \item $M(\ce{Zn})=\qty{65,4}{\gram\per\mole}$.
        \end{itemize}
    \end{donnees}
    On laisse fonctionner la pile (Zinc/Cuivre) de la figure ci-contre pendant
    une durée $\Delta t$ suffisamment longue jusqu'à ce que elle ne débite plus.
    
    L'équation bilan du fonctionnement de cette pile est~:
    \[
        \ce{Zn_{(s)} + Cu^{2+}_{(aq)} -> Zn^{2+}_{(aq)} + Cu_{(s)}}
    \]
\end{minipage}
\hfill
\begin{minipage}{0.34\linewidth}
    \begin{figure}[H]
        \centering
        \begin{tikzpicture}
            \newdimen\bh \bh=1.2cm
            \foreach \x/\c [count=\i] in {0/capri,3.5/lavendergray}{%
            \node[inner sep=0pt,minimum width=2cm,
                minimum height=\bh] (b\i) at (\x,0) {};
            \fill[\c,opacity=0.5,rounded corners]
                (b\i.north west) |- (b\i.south east) [sharp corners] --
                (b\i.north east) -- cycle;
            \draw[very thick,rounded corners,shorten <=-5mm,shorten >=-5mm]
                (b\i.north west) |- (b\i.south east) -- (b\i.north east);
            }
            % pont salin
            \coordinate (pb) at
                ([xshift=-5mm]$(b1.south east)!.4!(b1.north east)$);
            \coordinate (pe) at
                ([xshift=5mm]$(b2.south west)!.4!(b2.north west)$);
            \draw[double,thick,double distance=2pt,rounded corners]
                (pb) -- ++(0,1.5) coordinate(tmp) -- node[inner sep=2pt] (ps) {}
                (tmp -| pe) --(pe);
            \node[above,align=center] (pss) at (ps) {pont salin};
            % électrodes
            \node[inner sep=0pt,minimum width=4mm,minimum height=2cm,
                draw,thick,left color=white,right color=copper,
                anchor=south] (e1) at ([xshift=-1cm]pb) {};
            \node[inner sep=0pt,minimum width=4mm,minimum height=2cm,
                draw,thick,left color=white,right color=gray,
                anchor=south] (e2) at ([xshift=1cm]pe) {};
            % Text
            \draw[{Circle[scale=0.7]}-,shorten <=-5mm] (b1.south) -- ++(0,-0.3)
                node[below] {$\ce{Cu^{2+}_{(aq)} + SO$^{2-}_{4(aq)}$}$};
            \draw[{Circle[scale=0.7]}-,shorten <=-5mm] (b2.south) -- ++(0,-0.3)
                node[below] {$\ce{Zn^{2+}_{(aq)} + SO$^{2-}_{4(aq)}$}$};
            \node[below left,align=center] at (e1.north west) {$\ce{Cu}$};
            \node[below right,align=center] at (e2.north east) {$\ce{Zn}$};
            % Pôles
            %\node[circle,inner sep=4pt,draw,thick] (ca) at
            %    ([shift=(145:0.5)]e1.north west) {};
            %\draw[thick] (ca.east) -- (ca.west) (ca.north) -- (ca.south);
            %\node[circle,inner sep=4pt,draw,thick] (an) at
            %    ([shift=(35:0.5)]e2.north east) {};
            %\draw[thick] (an.east) -- (an.west);
            % Fils
            \draw (e1.north) to[short,i=$I$] ++(0,1.5)
                to[R,l=$R$] ++(2.4,0) coordinate(tmp)
                to[rmeter,t=A] (tmp -| e2.north) to[short] (e2.north);
        \end{tikzpicture}
    \end{figure}
\end{minipage}

\begin{enumerate}
    \item Donner le schéma conventionnel de cette pile.
    \item Montrer que la quantité de matière du cuivre déposé est
          $n(\ce{Cu})=\qty{5e-2}{\mole}$.
    \item Déterminer la valeur de la durée $\Delta t$ du fonctionnement de la
          pile sachant qu'elle délivre un courant d'intensité constante
          $I=\qty{100}{\mA}$.
\end{enumerate}

